\section{Perturbative solution of the Einstein--fluid equations}
\label{sec:einstein}

This section summarises the solution of $E_{\mu\nu}=G_{\mu\nu}[\eta^2\hat g]-8\pi G\,T_{\mu\nu}=0$ through
$O(\eps^3)$ in the parametrisation of Sec.~\ref{sec:setup}, obtained with a dedicated truncated-polynomial
engine (code \code{cmb4/code/C1\_einstein/}, App.~\ref{app:ledger}, C1-$\ast$).  The Christoffel symbols,
Ricci and Einstein tensors are computed directly from $g=\eta^2\hat g$ as series in $\eps$; the background
part of every component vanishes identically and the contracted Bianchi identity
$g^{\mu\alpha}\nabla_\alpha E_{\mu\nu}+g_{\nu\beta}\nabla_\mu T^{\mu\beta}=0$ holds identically through
$O(\eps^3)$.  Table~\ref{tab:objects} collects the named objects used in the rest of the paper; each is defined
once in the text and the long expressions are in App.~\ref{app:kernels}.

\begin{table}[t]
\centering\small
\begin{tabular}{llp{0.6\textwidth}}
\hline
object & defined in & meaning\\ \hline
$x=k\eta/\sqrt3$ & \eqref{eq:x} & dimensionless sound-horizon variable\\
$\TP,T_\delta,T_{\delta_T},T_V$ & \eqref{eq:transfer} & first-order transfer functions ($X_1=\zeta_1T_X$)\\
$\mathcal L_\Phi,\mathcal L_h,\mathcal L_\omega$ & \eqref{eq:operators} & universal linear operators (same at every order)\\
$S^{(n)}_{\mu\nu}$; $S_{\rm TL},S_{0}^{L},S^{T}_{0i},S^{V}_{ij},S^{TT}_{ij},S_\Phi$ & \eqref{eq:split}, \eqref{eq:sources} & order-$n$ sources and their SVT projections\\
$K_X(\vq,\vp;\eta)$, $\vp=\vk-\vq$ & \eqref{eq:kernel_def} & second-order kernel of $X$ per $\zeta_{\vq}\zeta_{\vp}$\\
$K_Q,K_{S_{00}},K_{S_{0z}},K_{S_\Phi},K_\omega,K_{v^T},K_h,K^0_{\Phi_2}$ & \eqref{eq:second_kernels}, App.~\ref{app:kernels} & the explicit second-order kernels\\
$u_1,u_2,\hat g,G_\Phi,G_h$ & \eqref{eq:greens} & homogeneous solutions and Green's functions\\
$\mathcal K_{\Phi_2},\mathcal K_h$ & \eqref{eq:declared2} & declared time integrals for $\Phi_2$, $h_2$\\
$X_{3c}(k,\eta)$, $K_{3c,X}(k,q,\mu;\eta)$ & \eqref{eq:contracted} & $q$-contracted third-order fields and kernels\\
$\mathrm{tad}_X(q,\eta)$ & \eqref{eq:tad} & mean-field (tadpole) functions $\langle X_2\rangle$\\
$C_3(k;q,\mu)$ & \eqref{eq:C3} & third-order superhorizon amplitude of $\Phi_{3c}$\\
\hline
\end{tabular}
\caption{Named objects of the perturbative solution.}
\label{tab:objects}
\end{table}

\subsection{First order}
\label{sec:first}
In the frame $\vk=k\hat{\mathbf z}$ the first-order equations are
\begin{align}
E^{(1)}_{ij}\ (i\neq j,\ \text{traceless}):&\quad -k^2(\Phi_1-\Psi_1)=0\ \Rightarrow\ \Psi_1=\Phi_1,\nonumber\\
E^{(1)}_{0i}:&\quad 2ik_i(\Phi_1'+\Psi_1/\eta)+4v_{1,i}/\eta^2=0\ \Rightarrow\ v^T_{1}=0,\quad
v_{1,i}=\partial_iV_1,\ V_1=-\tfrac{\eta^2}2(\Phi_1'+\Psi_1/\eta),\nonumber\\
E^{(1)}_{00}:&\quad \delta_1=-\tfrac23\eta^2k^2\Phi_1-2\eta\Phi_1'-2\Psi_1,\nonumber\\
\text{trace}:&\quad \Phi_1''+\tfrac4\eta\Phi_1'+\tfrac{k^2}3\Phi_1=0 .
\label{eq:first_eqs}
\end{align}
The velocity is irrotational and the conservation equations hold identically on this solution
(App.~\ref{app:ledger}, C1-3).  With
\begin{equation}
x\equiv k\eta/\sqrt3
\label{eq:x}
\end{equation}
the regular solution, normalised by $\zeta_1=\lim_{\eta\to0}(\delta_{T,1}-\Phi_1)$, is $X_1=\zeta_1T_X(x)$ with
\begin{align}
\TP(x)&=-\frac{2(\sin x-x\cos x)}{x^3}\ \xrightarrow{x\to0}\ -\tfrac23,\qquad \Psi_1=\Phi_1,\nonumber\\
T_\delta(x)&=4T_{\delta_T}(x)=-4\cos x+\frac{8\sin x}{x}+\frac{8\cos x}{x^2}-\frac{8\sin x}{x^3}\ \to\ \tfrac43,\nonumber\\
T_V(x)&=\eta\,\frac{x^2\sin x+2x\cos x-2\sin x}{x^3}\ \to\ \tfrac\eta3\qquad(v_{1,i}=\partial_iV_1,\ V_{1}=\zeta_1T_V).
\label{eq:transfer}
\end{align}
The second solution $(\cos x+x\sin x)/x^3\sim\eta^{-3}$ is discarded by regularity.  Note $\delta_{T,1}\to\zeta_1/3$,
so $\zeta_1=\delta_{T,1}-\Phi_1=\tfrac13\zeta_1+\tfrac23\zeta_1$ as it should.  We write
$T\equiv\TP$, $\dot T\equiv\partial_\eta\TP$, and $T_q\equiv\TP(q\eta/\sqrt3)$ etc.\ in the kernels.

\subsection{Second order: operator structure}
\label{sec:second}
At order $n\ge2$ every component of $E^{(n)}_{\mu\nu}=n!\,[\eps^n]E_{\mu\nu}$ splits as
\begin{equation}
E^{(n)}_{\mu\nu}=\mathrm{Lin}_{\mu\nu}[X_n]+S^{(n)}_{\mu\nu},
\label{eq:split}
\end{equation}
where $\mathrm{Lin}$ is the \emph{same} linear operator as in \eqref{eq:first_eqs} and $S^{(n)}_{\mu\nu}$ is built from
lower-order fields ($S^{(2)}$ quadratic in $X_1$; $S^{(3)}$ contains $X_1X_2$ and $X_1^3$ terms; position-space
expressions with 6--20 resp.\ 71--127 terms per component are in \code{position\_sources.tex}).  Decomposing
$S^{(n)}_{\mu\nu}$ into scalar, vector and tensor parts (frame $\vk=k\hat{\mathbf z}$, $\vq$ in the $x$--$z$ plane),
\begin{equation}
S_{\rm TL}\equiv S_{zz}-\tfrac12(S_{xx}+S_{yy}),\quad S_0^L\equiv S_{0z},\quad S^T_{0i}\equiv(S_{0x},S_{0y}),\quad
S^V\equiv S_{xz},\quad S^{TT}\equiv S_{xx}-S_{yy},\quad S_{\rm tr}\equiv S_{ii},
\label{eq:sources}
\end{equation}
the ten equations become (all verified by the engine):
\begin{align}
\text{traceless scalar:}&\quad Q_n\equiv\Phi_n-\Psi_n=S_{\rm TL}/k^2
\quad\big(=-\tfrac32\ilap(\partial_i\partial_j\ilap-\tfrac13\delta_{ij})S_{ij}\ \text{in position space}\big),\nonumber\\
00:&\quad \delta_n=\tfrac{\eta^2}3S_{00}-\tfrac23\eta^2k^2\Phi_n-2\eta\Phi_n'-2\Psi_n,\nonumber\\
0i\ \text{longitudinal:}&\quad v_{n,i}^L=\partial_iV_n,\qquad V_n=-\tfrac{\eta^2}4\Big[2(\Phi_n'+\Psi_n/\eta)+\frac{S_0^L}{ik}\Big],\nonumber\\
0i\ \text{transverse:}&\quad \tfrac12\lap\omega_{n,i}-\tfrac4{\eta^2}\big(\omega_{n,i}+v^T_{n,i}\big)=S^T_{0i}
\quad[8\pi Ga^2(\rho+P)=4/\eta^2],\nonumber\\
ij\ \text{vector:}&\quad \mathcal L_\omega\,\omega_{n}=\frac2{ik}S^V,\qquad \mathcal L_\omega\equiv\partial_\eta+\frac2\eta,\nonumber\\
ij\ \text{TT:}&\quad \mathcal L_h\,h_{n,ij}=-2\Lambda_{ij,kl}(\khat)S_{kl},\qquad \mathcal L_h\equiv\partial_\eta^2+\frac2\eta\partial_\eta+k^2,\nonumber\\
\text{trace:}&\quad \mathcal L_\Phi\,\Phi_n=S_\Phi\equiv-\frac{S_{\rm tr}-S_{00}}6+\frac{Q_n'}\eta-\frac{k^2}3Q_n,\qquad
\mathcal L_\Phi\equiv\partial_\eta^2+\frac4\eta\partial_\eta+\frac{k^2}3 ,
\label{eq:operators}
\end{align}
where $\Lambda_{ij,kl}=P_{ik}P_{jl}-\tfrac12P_{ij}P_{kl}$, $P_{ij}=\delta_{ij}-\hat k_i\hat k_j$, and in the trace
equation $\delta_n$ and $\Psi_n$ have been eliminated with the 00 and traceless equations.  Thus at every
order the scalar sector reduces to \emph{one} forced oscillator $\mathcal L_\Phi\Phi_n=S_\Phi$, after which
$\Psi_n,\delta_n,V_n$ are algebraic; the vector sector to a first-order equation; the tensor sector to a
forced wave equation.

\paragraph{Vector sector (relativistic Kelvin theorem).}
Vorticity-freedom of $v^i$ is \emph{not} implied at second order.  What the equations imply (verified: the
algebraic $\omega_2$ so obtained satisfies the $ij$-vector equation identically) is the circulation theorem
for the isentropic radiation fluid, whose enthalpy per particle is $\propto T$:
\begin{equation}
\partial_{[i}\big[(1+\delta_T)\,\hat u_{j]}\big]=0 .
\label{eq:kelvin}
\end{equation}
Consequently $\omega_2$ is \emph{algebraic} in the first-order fields and $v^T_2$ follows from the transverse
$0i$ constraint; no vector Green's function is needed (the vector ODE has $G_\omega=(\eta'/\eta)^2$ and gives the
same result).

\paragraph{Kernels.}
Second-order fields are written as
\begin{equation}
X_{2,\vk}(\eta)=\int\!\frac{\dd^3q}{(2\pi)^3}\,K_X(\vq,\vp;\eta)\,\zeta_{\vq}\zeta_{\vp},\qquad \vp\equiv\vk-\vq,
\qquad \mu\equiv\khat\cdot\qhat=\frac{k^2+q^2-p^2}{2kq},
\label{eq:kernel_def}
\end{equation}
and all scalar kernels are given in the \emph{invariant form} $K_X(k,q,p;\eta)$ (functions of the three
magnitudes), equivalently in the $(k,q,\mu)$ form.  With $T_q=\TP(q\eta/\sqrt3)$, $\dot T_q=\partial_\eta T_q$,
the second-order kernels obtained from \eqref{eq:operators} with $S^{(2)}$ are (full forms in App.~\ref{app:kernels}):
\begin{align}
K_Q&=T_qT_p\Big[-4+\frac{6\mu q}k+\frac{3q^2}{k^2}-\frac{9\mu^2q^2}{k^2}\Big]
+\Big[\eta(T_p\dot T_q+T_q\dot T_p)+\eta^2\dot T_q\dot T_p\Big]\Big[\frac{2\mu q}k+\frac{q^2}{k^2}-\frac{3\mu^2q^2}{k^2}\Big],\nonumber\\
K_\omega&=\frac{2i}{k^2}\Big[\eta\,(q^2-p^2)\,T_qT_p+\eta^2\big(q^2T_q\dot T_p-p^2T_p\dot T_q\big)\Big]
\quad(\omega_{2,i}=\int\!P_{ij}(\khat)q_j\,K_\omega\,\zeta\zeta;\ \text{antisymmetric in }q\leftrightarrow p),\nonumber\\
K_h&=4\Big[3T_qT_p+\eta(\dot T_qT_p+T_q\dot T_p)+\eta^2\dot T_q\dot T_p\Big]
=4\Big[2T_qT_p+(\eta\dot T+T)_q(\eta\dot T+T)_p\Big],\nonumber\\
K_{S_{00}},\ K_{S_{0z}},\ K_{S_\Phi},\ K_{v^T}&:\ \text{App.~\ref{app:kernels}};\qquad
v^T_{2,i}=\int P_{ij}q_jK_{v^T}\zeta\zeta,\qquad K_{v^T}=-\Big(\frac{\eta^2k^2}8+1\Big)K_\omega-\frac{\eta^2}4K_{S^T_{0i}} .
\label{eq:second_kernels}
\end{align}
The tensor source $K_h$ is identical to the radiation-era kernel of Kohri--Terada \cite{KT} (their eq.~(8)) and
Baumann et al.\ \cite{Baumann} (their (11)--(12)) with $w=\tfrac13$, $\Phi=\Psi$ (App.~\ref{app:ledger}, C1-9).
$K_{S_\Phi}$ is finite as $\eta\to0$ (no $\eta^{-2}$ terms), so the particular integral of the trace equation
vanishes at $\eta=0$; $\lim_{\eta\to0}\eta^2K_{S_{00}}=\tfrac{32}3$ and $K_Q(\eta\to0)$ is \eqref{eq:KQ0}.
The vector and tensor sources vanish as $\vp\to-\vq$ ($k\to0$) because of the projectors, which is what makes
the $K\to0$ corner of the $(22)$ loop integrable (Sec.~\ref{sec:monodef}).

\paragraph{Green's functions and the regular solution.}
With $x=k\eta/\sqrt3$, $x'=k\eta'/\sqrt3$,
\begin{align}
u_1(x)&=\frac{\sin x-x\cos x}{x^3},\qquad u_2(x)=\frac{\cos x+x\sin x}{x^3},\qquad
\hat g(x)\equiv3u_1(x)=-\tfrac32\TP(x)\ \to1,\qquad W\equiv u_1u_2'-u_1'u_2=-\frac{3\sqrt3}{\eta^4k^3},\nonumber\\
G_\Phi(k;\eta,\eta')&=\frac{u_1(x')u_2(x)-u_1(x)u_2(x')}{W(\eta')}
=\frac{\sqrt3}{k}\,\frac{x'}{x^3}\Big[(1+xx')\sin(x-x')-(x-x')\cos(x-x')\Big]
\ \xrightarrow{k\to0}\ \frac{\eta'}3\Big(1-\frac{{\eta'}^3}{\eta^3}\Big),\nonumber\\
G_h(k;\eta,\eta')&=\frac{\eta'}{\eta}\,\frac{\sin k(\eta-\eta')}{k},
\label{eq:greens}
\end{align}
satisfying $\mathcal L G=0$, $G(\eta,\eta)=0$, $\partial_\eta G|_{\eta'=\eta}=1$.  The regular second-order
scalar and tensor fields are then
\begin{align}
\Phi_{2,\vk}(\eta)&=\int\!\frac{\dd^3q}{(2\pi)^3}\zeta_{\vq}\zeta_{\vp}\,\mathcal K_{\Phi_2}(k,q,p;\eta),\qquad
\mathcal K_{\Phi_2}\equiv\int_{0}^{\eta}\!\dd\eta'\,G_\Phi(k;\eta,\eta')K_{S_\Phi}(k,q,p;\eta')+K^0_{\Phi_2}(k,q,p)\,\hat g(k\eta),\nonumber\\
h_{2,ij,\vk}(\eta)&=\Lambda_{ij,kl}(\khat)\int\!\frac{\dd^3q}{(2\pi)^3}q_kq_l\,\zeta_{\vq}\zeta_{\vp}\,\mathcal K_h(k,q,p;\eta),\qquad
\mathcal K_h\equiv\int_0^\eta\!\dd\eta'\,G_h(k;\eta,\eta')K_h(k,q,p;\eta'),
\label{eq:declared2}
\end{align}
with $K^0_{\Phi_2}$ the initial amplitude \eqref{eq:KPhi20} fixed by Gaussianity, and
$\Psi_2=\Phi_2-Q_2$, $\delta_2$, $V_2$ algebraic from \eqref{eq:operators}.  The lower limit $0$ may be replaced
by $\eta_i$ together with the homogeneous solution matched to the regular data at $\eta_i$ (Sec.~\ref{sec:regular}).
The time integrals $\mathcal K_{\Phi_2}$, $\mathcal K_h$ are \unevaluated{$\mathcal K_{\Phi_2}$, $\mathcal K_h$;
closed forms in terms of $\Si$, $\Ci$ of the Kohri--Terada type are being derived separately and will be quoted in
Sec.~\ref{sec:assembly}}.

\paragraph{Checks at second order.}
Bianchi propagation of the constraints (C1-5); exact dilation family \eqref{eq:dilation} (C1-6); tensor source
$=$ Kohri--Terada (C1-9); superhorizon $\Phi_2,\Psi_2$ $=$ Pitrou--Uzan--Bernardeau (C1-10); vorticity structure
\eqref{eq:kelvin} (C1-7).  An independent referee re-derivation (agent V3) is in progress and will be reported in
Sec.~\ref{sec:checks}.

\subsection{Third order: contracted kernels}
\label{sec:third}
The $(13)$ spectrum needs only the correlation of a third-order field with $\zeta_1$, i.e.\ the
\emph{$q$-contracted} third-order field,
\begin{equation}
\big\langle X_{3,\vk}\,\zeta_{\vk'}\big\rangle=(2\pi)^3\delta^3(\vk+\vk')\,P(k)\,X_{3c}(k,\eta),\qquad
X_{3c}(k,\eta)=\int\!\frac{\dd^3q}{(2\pi)^3}P(q)\,K_{3c,X}(k,q,\mu;\eta),
\label{eq:contracted}
\end{equation}
where the kernel sums the three Wick pairings (it is \emph{not} symmetrised under $\vq\to-\vq$, so only
$p=|\vk-\vq|$ appears).  The contracted source $K_{3c,S}$ is built from $S^{(3)}$ by the rules
\begin{equation}
A_1B_1C_1\ \to\ A_qB_{-q}C_k+A_qC_{-q}B_k+B_qC_{-q}A_k,\qquad
A_1B_2\ \to\ 2A_q\,K_B(-\vq,\vk)+A_k\,\langle B_2\rangle(\eta),
\label{eq:contraction_rules}
\end{equation}
in which $Q_2,\delta_2,V_2,v^T_2,\omega_2$ enter through their explicit kernels and $\Phi_2$, $h_2$ through the
declared integrals $\mathcal K_{\Phi_2}(a,b,c;\eta)$, $\mathcal K_h(a,b,c;\eta)$ with
$(a,b,c)=(|\vq_1|,|\vq_2|,|\vq_1+\vq_2|)$ kept as symbols.  The mean-field terms
\begin{equation}
\langle X_2\rangle(\eta)=\int\!\frac{\dd^3q}{(2\pi)^3}P(q)\,\mathrm{tad}_X(q,\eta),\qquad
\mathrm{tad}_X(q,\eta)\equiv\big\langle K_X(\vq,-\vq;\eta)\big\rangle_{\rm ang},
\label{eq:tad}
\end{equation}
vanish for vectors, tensors and velocities by isotropy; for the scalars $Q_2,S_{00},S_\Phi,\Phi_2,\delta_2$ the
angle-averaged $\vk\to0$ limits are given in closed form in App.~\ref{app:kernels}
(e.g.\ $\mathrm{tad}_{\Phi_2}(q,\eta)=-\tfrac89+\int_0^\eta\frac{\eta'}3(1-\frac{{\eta'}^3}{\eta^3})\mathrm{tad}_{S_\Phi}(q,\eta')\dd\eta'
=-\tfrac89-\tfrac{28}{405}q^2\eta^2+\ldots$), with the zero-mode caveat of Sec.~\ref{sec:tadpole}.

The contracted third-order equations have the same operators \eqref{eq:operators}:
\begin{align}
Q_{3c}&=K_{3c,Q},\qquad
\delta_{3c}=\tfrac{\eta^2}3K_{3c,S_{00}}-\tfrac23\eta^2k^2\Phi_{3c}-2\eta\Phi_{3c}'-2\Psi_{3c},\qquad
V_{3c}=-\tfrac{\eta^2}4\Big[2(\Phi_{3c}'+\Psi_{3c}/\eta)+\frac{K_{3c,S_{0z}}}{ik}\Big],\nonumber\\
\mathcal L_\Phi\Phi_{3c}&=K_{3c,S_\Phi},\qquad
\Phi_{3c}(k,\eta)=\int\!\frac{\dd^3q}{(2\pi)^3}P(q)\Big[\int_0^\eta\!\dd\eta'\,G_\Phi(k;\eta,\eta')K_{3c,S_\Phi}(k,q,\mu;\eta')+C_3(k;q,\mu)\,\hat g(k\eta)\Big],
\label{eq:third_eqs}
\end{align}
with $C_3$ the amplitude \eqref{eq:C3} fixed by $\zeta^{(3)}(0)=0$ and $K_{3c,S_\Phi}$ finite at $\eta\to0$.
The contracted vector and tensor sources ($S_{0x},S_{xz},S_{xx}-S_{yy}$) are produced as well but integrate to
zero over the $\vq$ angles by isotropy ($S_{0y},S_{yz},S_{xy}$ vanish identically).  The kernels
$K_{3c,X}$ have 157--830 terms and are not reproduced; they are stored in machine-readable form
(\code{third\_order.pkl}, loader \code{kernels\_api.load()['third']}, \LaTeX\ dump \code{third\_order.tex}).
The $\eta'$ integral in \eqref{eq:third_eqs} is \unevaluated{$\int_0^\eta G_\Phi K_{3c,S_\Phi}$, which itself
contains $\mathcal K_{\Phi_2}$ and $\mathcal K_h$ at the lower-order momenta}.
The un-contracted kernel $K_3(\vk_1,\vk_2,\vk_3)$ is not needed for $C^{(13)}_\ell$ and has not been produced.

\paragraph{Summary of what is settled.}
First order: closed forms \eqref{eq:transfer}.  Second order: all kernels in closed form except the two
time integrals $\mathcal K_{\Phi_2}$, $\mathcal K_h$, which are declared.  Third order: contracted sources and
amplitude $C_3$ in closed form; $\Phi_{3c}$ declared as one further nested time integral.  The zero-mode
convention for the tadpoles is a declared convention whose irrelevance for the observable is to be shown in
Sec.~\ref{sec:checks}.
